\chapter{Data Methods}

\section{Collection}

In the case of this project, there are two facets of required data:
\begin{enumerate}
  \item Open-Source Software Proliferation Data
  \item Job Postings Data
\end{enumerate}

\subsection{Open-Source Software Data}
Thankfully, the process of collecting data related to open-source software is fairly straightforward.
As the name suggests, the OSS projects are intended to be \emph{open} and accessible.
This necessitates the usage of some online platform for hosting and maintaining the codebase, with strong version-control systems to safeguard a maintainable codebase. \gh{} is functionally the default choice for hosting both open-source (and closed-source) codebases for a number of reasons, and is therefore the most relevant source for the required data of this project.

This project seeks data that more or less speaks to how \emph{popular} a given \ghrepo{} is, and the prime candidate for this is the \emph{\gh{} Stars} mechanism. Users will "star" a \repo{} in the same manner that internet users "bookmark" a website, behaving as a shortcut. This is a grossly simplfied view of how \ghrepo{} are interacted with, however - most \ghrepo{} have no stars, and most users have "starred" very few \repos{} in return. The star is instead a useful proxy for measuring interest from a development point of view - if someone is contributing regularly to a \repo{}, they will star it. In this sense, stars are a useful measure of how a \repo{} is gaining interest from developers (and in a lesser sense, users), which in turn represents a growth in the utility the \repo{} represents.

\cite{__GitHubDaily} is an open-source project that provides data on the historic star-counts of \ghrepo{}, marketed as a tool to "provide insights into a \repo{}'s trendiness". Within the ML/AI framework that I delineated, the relevant \repos{}' star counts were found and downloaded thusly.

\subsection{Job Posting Data}
As is generally true of commercially valuable data, it can prove difficult to find open data that properly satisfies the desired criteria. Generally, it is a relatively simple task to build legally compliant scrapers that programatically access and store the data from job boards online. However, finding stores of historical data presents a major challenge, because only \emph{active} job postings are present on job boards. This means that the scraper must be active throughout and is not useful for collecting data in retrospect.

An alternative solution is to therefore find and access a pre-prepared dataset. This, unfortunately, proved challenging in the sense that:
\begin{itemize}
  \item Large, high-quality datasets are presented as products, and are thus expensive
  \item Small, low-quality datasets often lack essential features, such as "job description" and even "date"
\end{itemize}

The solution to these problems was to utilize \emph{several} datasets that occupy different time-regions, made available on Kaggle. An overview of the characteristics for each dataset is included in the Appendix, as per Table~\ref{table:job_data_overview}.

At this rate,